\documentclass[10pt,a4paper]{amsart}
\usepackage[margin=19mm]{geometry}
\usepackage{amsmath,amssymb,mathtools}
\usepackage[scr=rsfs,cal=euler]{mathalfa}
\usepackage{xfrac}
\usepackage[unicode,hidelinks,bookmarks=false]{hyperref}
\newcommand{\ie}{i.e.\ }
\newcommand{\cf}{cf.\ }
\newcommand{\resp}{respectively\ }
\newcommand{\Subsection}{\subsection}
\providecommand{\nobreakdash}{\nobreak\hbox{-}}
\setlength{\emergencystretch}{2em}
\multlinegap0pt
\usepackage{bm}
\usepackage{bbm}
\usepackage{dsfont}
\usepackage{xspace}
\usepackage{cancel}
\usepackage{mathtools}
\usepackage{amsthm}
\makeatletter
\@ifundefined{theorem}{\newtheorem{theorem}{Theorem}}{}
\@ifundefined{lemma}{\newtheorem{lemma}[theorem]{Lemma}}{}
\@ifundefined{proposition}{\newtheorem{proposition}[theorem]{Proposition}}{}
\makeatother
\title[A reprise of the NTV conjecture for the Hilbert transform]{A reprise of the NTV conjecture for the Hilbert transform}
\author{Eric T. Sawyer}
\date{Source: arXiv:2302.13920v13 (CC BY 4.0)}
\begin{document}
\maketitle

\noindent\emph{Source and licence.} A reprise of the NTV conjecture for the Hilbert transform. Version arXiv:2302.13920v13 (CC BY 4.0), distributed under the Creative Commons Attribution 4.0 International licence, as recorded in the arXiv licence metadata for this exact version and shown on the abstract page https://arxiv.org/abs/2302.13920v13. Full rights notice: licenses/arxiv-2302.13920-CC-BY-4.0.txt.\par\smallskip

\noindent\textit{Editorial scope.} Counted unit: the Proposition whose source label is \texttt{strongly adapted'} in the article (source0.tex lines 949--962, source label \texttt{strongly adapted'}), delivered together with the article's own complete original proof of it (source0.tex lines 964--1140, one \texttt{proof} environment of 177 lines). No mathematics is written, repaired or re-derived: statement and proof are the article's own text line for line. Required context retained uncounted: Energy Lemma (lines 224--243). Every definition, display and label of those lines is retained unchanged. Omitted from this package and not redistributed: 1--223, 244--948, 963--963, 1141--2975. Nothing inside a retained excerpt is omitted or altered: every retained body is delivered exactly as the article's lines are written, so no omission receipt of any kind accompanies this package. Citations: the retained excerpts carry no citation command at all, so no bibliography entry of the article is transcribed and no reference is redistributed. Reference adaptation: none was needed and none was made; every reference inside the delivered unit resolves inside it, so no unresolved reference or citation is produced. Printed slips kept literal: no printed word, symbol, hypothesis or number of the counted statement or of its proof was altered. Layout and class: the article's own class is \texttt{amsart} with packages including amssymb, amsfonts, geometry, amsmath, graphicx; the wrapper is the lane's pinned \texttt{amsart} editorial wrapper at 10pt on A4, which restates the theorem-like environments the retained text needs and adds the article's own macro definitions that the retained text needs (none), with extra wrapper packages (bm, bbm, dsfont, xspace, cancel, mathtools). The delivered numbering is the unit's own. Assets: no retained span contains a figure environment, a graphics-inclusion command, a PDF-page inclusion, a table or a TikZ picture; no image file is copied into this package, the package declares no asset dependency and no image file is redistributed with it. Licence: version arXiv:2302.13920v13 of this article is distributed under the Creative Commons Attribution 4.0 International licence, as recorded in the arXiv licence metadata for this exact version and shown on the abstract page https://arxiv.org/abs/2302.13920v13. A full rights notice is delivered at licenses/arxiv-2302.13920-CC-BY-4.0.txt. Characters: every retained excerpt is pure ASCII, so the delivered engine, XeLaTeX with its default Latin Modern text font, reports no missing character.
\begin{lemma}
[Monotonicity and Energy Lemma]\label{Energy Lemma}Suppose that $J\subset
I^{\prime}\subset2I^{\prime}\subset F$ are dyadic intervals. Then
\begin{align*}
\left\vert \left\langle H_{\sigma}\mathbf{1}_{F\setminus I^{\prime}%
},\bigtriangleup_{J}^{\omega}g\right\rangle _{\omega}\right\vert ^{2}  &
\lesssim\left(  \frac{\mathrm{P}\left(  J,\mathbf{1}_{F\setminus I^{\prime}%
}\sigma\right)  }{\ell\left(  J\right)  }\right)  ^{2}\left\Vert
\bigtriangleup_{J}^{\omega}x\right\Vert _{L^{2}\left(  \omega\right)  }%
^{2}\left\Vert \bigtriangleup_{J}^{\omega}g\right\Vert _{L^{2}\left(
\omega\right)  }^{2}\ ,\\
\left\vert \left\langle H_{\sigma}\mathbf{1}_{F\setminus I^{\prime}}%
,\sum_{J^{\prime}\subset J}a_{J^{\prime}}h_{J^{\prime}}^{\omega}\right\rangle
_{\omega}\right\vert ^{2}  & \lesssim\mathrm{P}\left(  J,\mathbf{1}%
_{F\setminus I^{\prime}}\sigma\right)  ^{2}\left\vert J\right\vert _{\omega
}\left\Vert \sum_{J^{\prime}\subset J}a_{J^{\prime}}h_{J^{\prime}}^{\omega
}\right\Vert _{L^{2}\left(  \omega\right)  }^{2}.
\end{align*}

\end{lemma}

\begin{proposition}
[The Intertwining Proposition]\label{strongly adapted'}Suppose that
$\mathcal{F}$ is $\sigma$-Carleson. Then%
\[
\left\vert \sum_{F\in\mathcal{F}}\ \sum_{I:\ I\supsetneqq F}\ \left\langle
H_{\sigma}\left(  \mathbf{1}_{I_{F}}\bigtriangleup_{I}^{\sigma}f\right)
,\mathsf{P}_{\mathcal{C}_{\mathcal{F}}\left(  F\right)  }^{\omega
}g\right\rangle _{\omega}\right\vert \lesssim\left(  \mathfrak{F}\left(
\sigma,\omega\right)  +\mathfrak{T}_{H}\left(  \sigma,\omega\right)  \right)
\ \left\Vert f\right\Vert _{L^{2}\left(  \sigma\right)  }\left\Vert
g\right\Vert _{L^{2}\left(  \omega\right)  }.
\]

\end{proposition}

\begin{proof}
We write the left hand side of the display above as%
\[
\sum_{F\in\mathcal{F}}\ \sum_{I:\ I\supsetneqq F}\ \left\langle H_{\sigma
}\left(  \mathbf{1}_{I_{F}}\bigtriangleup_{I}^{\sigma}f\right)  ,g_{F}%
\right\rangle _{\omega}=\sum_{F\in\mathcal{F}}\ \left\langle H_{\sigma}\left(
\sum_{I:\ I\supsetneqq F}\mathbf{1}_{I_{F}}\bigtriangleup_{I}^{\sigma
}f\right)  ,g_{F}\right\rangle _{\omega}\equiv\sum_{F\in\mathcal{F}%
}\ \left\langle H_{\sigma}f_{F},g_{F}\right\rangle _{\omega}\ ,
\]
where%
\[
g_{F}=\mathsf{P}_{\mathcal{C}_{\mathcal{F}}\left(  F\right)  }^{\omega}%
g=\sum_{J\in\mathcal{C}_{\mathcal{F}}\left(  F\right)  }\bigtriangleup
_{J}^{\omega}g\text{ and }f_{F}\equiv\sum_{I:\ I\supsetneqq F}\mathbf{1}%
_{I_{F}}\bigtriangleup_{I}^{\sigma}f\ .
\]
Note that $g_{F}$ is supported in $F$, and that $f_{F}$ is constant on $F$.
The intervals $I$ occurring in this sum are linearly and consecutively ordered
by inclusion, along with the intervals $F^{\prime}\in\mathcal{F}$ that contain
$F$. More precisely, we can write%
\[
F\equiv F_{0}\subsetneqq F_{1}\subsetneqq F_{2}\subsetneqq...\subsetneqq
F_{n}\subsetneqq F_{n+1}\subsetneqq...F_{N}%
\]
where $F_{m}=\pi_{\mathcal{F}}^{m}F$ for all $m\geq1$. We can also write%
\[
F=F_{0}\subsetneqq I_{1}\subsetneqq I_{2}\subsetneqq...\subsetneqq
I_{k}\subsetneqq I_{k+1}\subsetneqq...\subsetneqq I_{K}=F_{N}%
\]
where $I_{k}=\pi_{\mathcal{D}}^{k}F$ for all $k\geq1$, and by convention we
set $I_{0}=F$. There is a (unique) subsequence $\left\{  k_{m}\right\}
_{m=1}^{N}$ such that%
\[
F_{m}=I_{k_{m}},\ \ \ \ \ 1\leq m\leq N.
\]


Recall that%
\[
f_{F}\left(  x\right)  =\sum_{k=1}^{\infty}\mathbf{1}_{\left(  I_{k}\right)
_{F}}\left(  x\right)  \bigtriangleup_{I_{k}}^{\sigma}f\left(  x\right)
=\sum_{k=1}^{\infty}\mathbf{1}_{I_{k}\setminus I_{k-1}}\left(  x\right)
\sum_{\ell=k+1}^{\infty}\bigtriangleup_{I_{\ell}}^{\sigma}f\left(  x\right)  .
\]
Assume now that $k_{m}\leq k<k_{m+1}$. Using a telescoping sum, we compute
that for
\[
x\in I_{k+1}\setminus I_{k}\subset F_{m+1}\setminus F_{m},
\]
we have
\[
\left\vert f_{F}\left(  x\right)  \right\vert =\left\vert \sum_{\ell
=k+2}^{\infty}\bigtriangleup_{I_{\ell}}^{\sigma}f\left(  x\right)  \right\vert
=\left\vert \mathbb{E}_{\theta I_{k+2}}^{\sigma}f-\mathbb{E}_{I_{K}}^{\sigma
}f\right\vert \lesssim\mathbb{E}_{F_{m+1}}^{\sigma}\left\vert f\right\vert \ .
\]
Now $f_{F}$ is constant on $F$ and%
\begin{align*}
\left\vert f_{F}\right\vert  & \leq\sum_{m=0}^{N}\left(  \mathbb{E}_{F_{m+1}%
}^{\sigma}\left\vert f\right\vert \right)  \ \mathbf{1}_{F_{m+1}\setminus
F_{m}}=\left(  \mathbb{E}_{F}^{\sigma}\left\vert f\right\vert \right)
\ \mathbf{1}_{F}+\sum_{m=0}^{N}\left(  \mathbb{E}_{\pi_{\mathcal{F}}^{m+1}%
F}^{\sigma}\left\vert f\right\vert \right)  \ \mathbf{1}_{\pi_{\mathcal{F}%
}^{m+1}F\setminus\pi_{\mathcal{F}}^{m}F}\\
& =\left(  \mathbb{E}_{F}^{\sigma}\left\vert f\right\vert \right)
\ \mathbf{1}_{F}+\sum_{F^{\prime}\in\mathcal{F}:\ F\subset F^{\prime}}\left(
\mathbb{E}_{\pi_{\mathcal{F}}F^{\prime}}^{\sigma}\left\vert f\right\vert
\right)  \ \mathbf{1}_{\pi_{\mathcal{F}}F^{\prime}\setminus F^{\prime}}\\
& \leq\alpha_{\mathcal{F}}\left(  F\right)  \ \mathbf{1}_{F}+\sum_{F^{\prime
}\in\mathcal{F}:\ F\subset F^{\prime}}\alpha_{\mathcal{F}}\left(
\pi_{\mathcal{F}}F^{\prime}\right)  \ \mathbf{1}_{\pi_{\mathcal{F}}F^{\prime
}\setminus F^{\prime}}\\
& \leq\alpha_{\mathcal{F}}\left(  F\right)  \ \mathbf{1}_{F}+\sum_{F^{\prime
}\in\mathcal{F}:\ F\subset F^{\prime}}\alpha_{\mathcal{F}}\left(
\pi_{\mathcal{F}}F^{\prime}\right)  \ \mathbf{1}_{\pi_{\mathcal{F}}F^{\prime}%
}\ \mathbf{1}_{F^{c}}\\
& =\alpha_{\mathcal{F}}\left(  F\right)  \ \mathbf{1}_{F}+\Phi\ \mathbf{1}%
_{F^{c}}\ ,\ \ \ \ \ \text{\ for all }F\in\mathcal{F},
\end{align*}
where%
\[
\Phi\equiv\sum_{F^{\prime\prime}\in\mathcal{F}}\ \alpha_{\mathcal{F}}\left(
F^{\prime\prime}\right)  \ \mathbf{1}_{F^{\prime\prime}}\ .
\]


Now we write%
\[
\sum_{F\in\mathcal{F}}\ \left\langle H_{\sigma}f_{F},g_{F}\right\rangle
_{\omega}=\sum_{F\in\mathcal{F}}\ \left\langle H_{\sigma}\left(
\mathbf{1}_{F}f_{F}\right)  ,g_{F}\right\rangle _{\omega}+\sum_{F\in
\mathcal{F}}\ \left\langle H_{\sigma}\left(  \mathbf{1}_{F^{c}}f_{F}\right)
,g_{F}\right\rangle _{\omega}\equiv I+II.
\]
Then local interval testing and quasi-orthogonality, together with the fact
that $f_{F}$ is a constant on $F$ bounded by $\alpha_{\mathcal{F}}\left(
F\right)  $, give%
\begin{align*}
& \left\vert I\right\vert =\left\vert \sum_{F\in\mathcal{F}}\int_{\mathbb{R}%
}\mathbf{1}_{F}\left(  x\right)  H_{\sigma}\left(  \mathbf{1}_{F}f_{F}\right)
\left(  x\right)  \ g_{F}\left(  x\right)  \ d\omega\left(  x\right)
\right\vert \\
& \leq\sqrt{\sum_{F\in\mathcal{F}}\int_{\mathbb{R}}\left\vert \alpha
_{\mathcal{F}}\left(  F\right)  \mathbf{1}_{F}\left(  x\right)  H_{\sigma
}\left(  \mathbf{1}_{F}\right)  \left(  x\right)  \right\vert ^{2}%
d\omega\left(  x\right)  }\ \sqrt{\sum_{F\in\mathcal{F}}\int_{\mathbb{R}%
}\left\vert g_{F}\left(  x\right)  \right\vert ^{2}d\omega\left(  x\right)
}\\
& \lesssim\mathfrak{T}_{H}^{\operatorname{loc}}\left(  \sigma,\omega\right)
\sqrt{\sum_{F\in\mathcal{F}}\int_{\mathbb{R}}\left\vert \alpha_{\mathcal{F}%
}\left(  F\right)  \right\vert ^{2}\left\vert F\right\vert _{\sigma}%
}\left\Vert g\right\Vert _{L^{2}\left(  \omega\right)  }\lesssim
\mathfrak{T}_{H}^{\operatorname{loc}}\left(  \sigma,\omega\right)  \left\Vert
f\right\Vert _{L^{2}\left(  \sigma\right)  }\left\Vert g\right\Vert
_{L^{2}\left(  \omega\right)  }.
\end{align*}


Now $\mathbf{1}_{F^{c}}f_{F}$ is supported outside $F$, and each $J$ in the
Haar support of $g_{F}=\mathsf{P}_{\mathcal{C}_{\mathcal{F}}\left(  F\right)
}^{\omega}g$ is either in $\mathcal{N}_{r}\left(  F\right)  =\left\{
J\in\mathcal{D}\left[  F\right]  :\ell\left(  J\right)  \geq2^{-r}\ell\left(
F\right)  \right\}  $, in which case the desired bound for term $I$ is
straightforward, or $J$ is $\left(  r,\varepsilon\right)  $-deeply embedded in
$F$, i.e. $J\subset_{r,\varepsilon}F$, and so $J\subset_{r,\varepsilon}W$ for
some $W\in\mathcal{M}_{\left(  r,\varepsilon\right)  -\operatorname*{deep}%
}\left(  F\right)  $. Thus with $\mathcal{C}_{\mathcal{F}}^{\flat}\left(
F\right)  \equiv\mathcal{C}_{\mathcal{F}}\left(  F\right)  \setminus
\mathcal{N}_{r}\left(  F\right)  $ we have%
\begin{align*}
\left\vert II\right\vert  & =\left\vert \sum_{F\in\mathcal{F}}\int
_{\mathbb{R}}H_{\sigma}\left(  \mathbf{1}_{F^{c}}f_{F}\right)  \left(
x\right)  \ \mathsf{P}_{\mathcal{C}_{\mathcal{F}}^{\flat}\left(  F\right)
}^{\omega}g\left(  x\right)  \ d\omega\left(  x\right)  \right\vert \\
& =\left\vert \sum_{F\in\mathcal{F}}\sum_{W\in\mathcal{M}_{\left(
r,\varepsilon\right)  -\operatorname*{deep}}\left(  F\right)  \cap
\mathcal{C}_{\mathcal{F}}\left(  F\right)  }\int_{\mathbb{R}}\mathsf{P}%
_{\mathcal{C}_{\mathcal{F}}^{\flat}\left(  F\right)  \cap\mathcal{D}\left[
W\right]  }^{\omega}H_{\sigma}\left(  \mathbf{1}_{F^{c}}f_{F}\right)  \left(
x\right)  \ \mathsf{P}_{\mathcal{C}_{\mathcal{F}}^{\flat}\left(  F\right)
\cap\mathcal{D}\left[  W\right]  }^{\omega}g\left(  x\right)  \ d\omega\left(
x\right)  \right\vert \\
& \lesssim\sqrt{\sum_{F\in\mathcal{F}}\sum_{W\in\mathcal{M}_{\left(
r,\varepsilon\right)  -\operatorname*{deep}}\left(  F\right)  \cap
\mathcal{C}_{\mathcal{F}}\left(  F\right)  }\int_{\mathbb{R}}\left\vert
\mathsf{P}_{\mathcal{C}_{\mathcal{F}}^{\flat}\left(  F\right)  \cap
\mathcal{D}\left[  W\right]  }^{\omega}H_{\sigma}\left(  \mathbf{1}_{F^{c}%
}f_{F}\right)  \left(  x\right)  \right\vert ^{2}d\omega\left(  x\right)  }\\
& \ \ \ \ \ \ \ \ \ \ \ \ \ \ \ \ \ \ \ \ \ \ \ \ \ \times\sqrt{\sum
_{F\in\mathcal{F}}\sum_{W\in\mathcal{M}_{\left(  r,\varepsilon\right)
-\operatorname*{deep}}\left(  F\right)  \cap\mathcal{C}_{\mathcal{F}}\left(
F\right)  }\int_{\mathbb{R}}\left\vert \mathsf{P}_{\mathcal{C}_{\mathcal{F}%
}^{\flat}\left(  F\right)  \cap\mathcal{D}\left[  W\right]  }^{\omega}g\left(
x\right)  \right\vert ^{2}d\omega\left(  x\right)  }.
\end{align*}
The second factor is bounded by $C\left\Vert g\right\Vert _{L^{2}\left(
\omega\right)  }$, and we use the second line in the Energy Lemma
\ref{Energy Lemma} on the first factor to bound it by,%
\begin{align}
& \sqrt{\sum_{F\in\mathcal{F}}\sum_{W\in\mathcal{M}_{\left(  r,\varepsilon
\right)  -\operatorname*{deep}}\left(  F\right)  \cap\mathcal{C}_{\mathcal{F}%
}\left(  F\right)  }\ \left(  \frac{\mathrm{P}\left(  W,\mathbf{1}_{F^{c}%
}f_{F}\sigma\right)  }{\ell\left(  W\right)  }\right)  ^{2}\int_{W}\left\vert
\mathsf{P}_{\mathcal{C}_{\mathcal{F}}\left(  F\right)  \cap\mathcal{D}\left[
W\right]  }^{\omega}x\right\vert ^{2}d\omega}\label{bound by}\\
& \leq\mathfrak{F}\left(  \sigma,\omega\right)  \lVert\Phi\rVert_{L^{p}\left(
\sigma\right)  }\lesssim\mathfrak{F}\left(  \sigma,\omega\right)  \lVert
f\rVert_{L^{2}\left(  \sigma\right)  },\nonumber
\end{align}
using the definition of the functional energy characteristic $\mathfrak{F}%
\left(  \sigma,\omega\right)  $, and the dyadic $\sigma$-maximal function
inequality $\left\Vert \Phi\right\Vert _{L^{p}\left(  \sigma\right)  }%
\lesssim\left\Vert M_{\sigma}f\right\Vert _{L^{p}\left(  \sigma\right)
}\lesssim\left\Vert f\right\Vert _{L^{p}\left(  \sigma\right)  }$. This
completes the proof of the Intertwining Proposition \ref{strongly adapted'}.
\end{proof}

\end{document}
